%% =====================================================================
%%  SUPPLEMENTARY MATERIAL
%%  "Efficient and practical solutions for the correlated storage
%%   location assignment problem in automated pick-and-pass warehouses"
%%  International Journal of Production Research
%%
%%  This file holds the calibration and sensitivity study of the
%%  clustering heuristic, referenced from Section 4.1 of the main
%%  article. It compiles standalone: it needs only IJPR_CSLAP.bib
%%  alongside it. Cross-references to the main article are written out
%%  in full, since the two documents compile separately.
%% =====================================================================
\PassOptionsToPackage{table}{xcolor}
\documentclass[11pt]{article}

\usepackage[a4paper,margin=2.5cm]{geometry}
\usepackage{amsmath,amssymb,amsfonts}
\usepackage{mathtools}
\usepackage{booktabs}
\usepackage{float}
\usepackage{xcolor}
\usepackage{tabularx}
\usepackage{multirow}
\usepackage{algorithmicx}
\usepackage{algorithm}
\usepackage{algpseudocode}
\usepackage{url}
\usepackage{hyperref}

\usepackage[natbibapa,nodoi]{apacite}
\setlength\bibhang{12pt}

% Number sections, tables and figures with an S prefix, as T&F expect
% for supplementary items.
\renewcommand{\thesection}{S\arabic{section}}
\renewcommand{\thetable}{S\arabic{table}}
\renewcommand{\thefigure}{S\arabic{figure}}
\renewcommand{\thealgorithm}{S\arabic{algorithm}}

% The main article uses \tbl{} and \tabnote{} from interact.cls; supply
% equivalents so the table bodies below are unchanged from the article.
\newcommand{\tbl}[1]{\caption{#1}}
\newcommand{\tabnote}[1]{\vspace{2pt}\footnotesize\raggedright #1\par}

\title{Supplementary material:\\
Calibration and sensitivity of the clustering heuristic}
\author{Ermal Belul \and Marwane Bouznif \and Dritan Nace \and Antoine Jouglet}
\date{}

\begin{document}
\maketitle

\noindent
This supplement evaluates the design parameters of the greedy clustering heuristic of
Section~4.1 of the main article across two distinct synthetic experimental setups. We report
two investigations: (1) determining whether the community size bound $\beta$ correlates
primarily with total warehouse catalogue size ($N$) or with station location capacity
($\zeta$), and (2) evaluating the sensitivity of solution quality to variations in each
heuristic threshold ($\phi, \gamma, r, \beta$).

To isolate these effects without confounding catalogue volume and station capacity (which are
tied in the main 29-instance benchmark of Section~5 of the main article), we generated an
expanded dataset of 80 synthetic instances across crossed combinations of $N \in \{500, 1000\}$
and $\zeta \in \{20, 25, 50, 100\}$, using the same common-itemset generator script as the
primary benchmark but with independent random seeds. This secondary dataset allows us to
systematically test whether $\beta$ scales with $N$ or $\zeta$, and to test threshold
robustness across diverse warehouse configurations.

\section{Community bound scaling across warehouse geometries}\label{sup:capsweep}
The three filters of Section~4.1 of the main article are indexed on the catalogue and the
community bound on the station. That asymmetry needs evidence, because the benchmark of
Section~5 cannot supply it: its generator ties the station count to the catalogue, so
$\zeta=100$ at every size from 500 SKUs upward and the ratio $\beta/\zeta$ is constant across
the whole family. Catalogue size and slot capacity are confounded there, and no result measured
on it can separate them.

We therefore built a second family in which the two are not confounded. It uses the same
generator as Section~5 of the main article (the common-itemset procedure of
\citet{zhang2019cie} at correlation degree $0.7$, with order size drawn from $U[5,15]$ and
proportional order volumes) and changes one thing: the number of stations is set independently
of the catalogue rather than as $\max(5,\lfloor N/100\rfloor)$. Two catalogue sizes,
$N\in\{500,1000\}$, are crossed with four station counts chosen so that the same slot
capacities $\zeta\in\{20,25,50,100\}$ occur at both sizes. Ten instances per cell give 80 in
all, generated on seeds $2001$ to $2010$, disjoint from the seeds of the benchmark. Station
labels are drawn from a separate random stream, so at a fixed seed the order structure is
identical across the four geometries and only the station configuration changes. Any difference
between cells is therefore attributable to geometry alone.

Each instance is solved at every $\beta$ on the grid
$\{2,3,5,8,10,12,15,18,20,25,30,40,50,75,100\}$ that falls at or below its own $\zeta$, which
gives 980 runs. All 980 returned a feasible layout. Table~\ref{tab:capsweep} reports the
minimising bound in each cell. Both instance families, with the generator settings and seeds
needed to reconstruct them exactly, accompany the article.

\begin{table}[H]
\centering
\tbl{Visit-minimising community bound $\beta^\ast$ across warehouse geometries, on 80 instances disjoint from the benchmark. The same $\zeta$ appears at both catalogue sizes.}
{\small\begin{tabular*}{\textwidth}{@{\extracolsep{\fill}}rrrrrr@{}}
\toprule
$\zeta$ & $N$ & $|S|$ & $\beta^\ast$ & $\beta^\ast/\zeta$ & $\Delta$ visits vs.\ $\beta=15$ \\
\midrule
20  & 500   & 25 & 12 & 0.60 & $-0.72$\% \\
20  & 1{,}000 & 50 & 12 & 0.60 & $-0.45$\% \\
25  & 500   & 20 & 15 & 0.60 & $0.00$\% \\
25  & 1{,}000 & 40 & 15 & 0.60 & $0.00$\% \\
50  & 500   & 10 & 30 & 0.60 & $-0.87$\% \\
50  & 1{,}000 & 20 & 30 & 0.60 & $-0.84$\%\textsuperscript{$\ast$} \\
100 & 500   & 5  & 40 & 0.40 & $-1.01$\%\textsuperscript{$\ast$} \\
100 & 1{,}000 & 10 & 40 & 0.40 & $-1.06$\%\textsuperscript{$\ast$} \\
\bottomrule
\end{tabular*}}
\tabnote{$\beta^\ast$ minimises mean visits in its cell over the swept values, which include 15. The last column is the change from 15 to $\beta^\ast$, so a negative entry means $\beta^\ast$ uses fewer visits and a zero entry means 15 is itself the minimiser, as it is in both $\zeta=25$ cells. \textsuperscript{$\ast$}Separates from 15 at $p<0.05$ under a Holm-corrected Wilcoxon signed-rank test over the ten instances of the cell. The unstarred non-zero cells fail that test on within-cell variance rather than on effect size: at $\zeta=50$, $N=500$ the point effect of $-0.87$\% exceeds that of the significant cell beneath it.}
\label{tab:capsweep}
\end{table}

At a fixed slot capacity $\beta^\ast$ is the same at both catalogue sizes in all four cells,
whereas at a fixed catalogue size it moves with $\zeta$ across the whole range. Over the eight
cells the rank correlation between $\beta^\ast$ and $\zeta$ is $0.962$, and between
$\beta^\ast$ and $N$ it is $0.000$. This confirms that the bound scales with the slot capacity
of a station and is insensitive to the length of the catalogue.

The dependence is monotone but not proportional. The minimising ratio $\beta^\ast/\zeta$ is
$0.60$ at $\zeta\le 50$ and falls to $0.40$ at $\zeta=100$, so a single coefficient cannot be
optimal everywhere. We adopt $0.4$ for two reasons. It is the minimising value at $\zeta=100$,
the capacity of every benchmark instance of 500 SKUs and above and of the sizes at which the
method is intended to be used. Where it is not the minimiser it errs low, which matters because
larger communities are harder to place within the workload constraint, as Section~6 of the main
article shows on a site whose stations are not interchangeable. The price of that choice is nil
at $\zeta=100$, where $0.4\zeta$ is the minimiser, and between $0.15$ and $1.03$\% of visits in
the six cells below it.

\section{Threshold sensitivity of the clustering heuristic}\label{sup:sensitivity}
Table~\ref{tab:sensitivity} rescales each of the four thresholds of Section~4.1 of the main
article from half to twice its nominal value, one at a time and then all four together, on the
twelve 50-SKU and ten 500-SKU instances. The 1,000- and 2,000-SKU instances are omitted because
four and three instances are too few for an informative interval, although the ordering of the
effects there agrees with what follows. Multipliers apply to the effective nominal value on the
instance, after the floor of 5 (for the three filters that value is set by $N$ and for the
community bound by $\zeta$). At 50 SKUs both frequency floors already sit on their clamp floor
of 2, and the community bound sits on its floor of 5, so only the upward settings are distinct
there. Entries are paired mean changes against the nominal setting on the same instance. The
procedure is deterministic: repeating the nominal setting under five hash seeds returns an
identical layout on every instance, so any difference is a real threshold effect and not
run-to-run variation. Reporting follows the conventions for computational experiments with
heuristics \citep{barr1995designing, hooker1995testing}.

The two frequency floors are inert, as rescaling either leaves the visit count unchanged and
the surviving edge set bit-identical at every size, because neither removes a product or pair
that the kept-edge ratio does not already remove. The ratio changes the edge set substantially
but the objective very little, separating from the nominal setting only at $\times2$. The
community bound is the one threshold whose response is not monotone. Halving it costs $0.63\%$
at 500 SKUs and doubling it costs $1.19\%$, while $\times0.75$ and $\times1.5$ are
indistinguishable from the nominal setting, while at 1,000 and 2,000 SKUs the pattern is the
same, with doubling costing $1.84\%$ and $2.27\%$. The exception is $\times1.5$ at 500 SKUs,
which recovers $0.45\%$, consistent with Section~\ref{sup:capsweep} above, where the minimising
ratio at $\zeta=100$ is $0.40$ but rises to $0.60$ at smaller capacities, so a bound half again
as large is not yet past the minimum on every instance. A default whose response rises on both
sides at the two larger sizes, and is flat rather than falling at the smaller one, is indexed on
the right quantity, whereas a default with a monotone response is not. Leaving all four at their
defaults is therefore supported here, not assumed.

\begin{table}[H]
\centering
\tbl{Sensitivity of the clustering heuristic to its four thresholds, on the 22 instances of 50 and 500 SKUs.}
{\small\setlength{\tabcolsep}{4pt}\begin{tabularx}{\textwidth}{@{}X c c@{}}
\toprule
Threshold setting & $\Delta$ visits, 50 SKUs & $\Delta$ visits, 500 SKUs \\
\midrule
Frequency floor $\phi$, $\times0.5$ to $\times2$\textsuperscript{a}     & $0.00$\% & $0.00$\% \\
Co-occurrence floor $\gamma$, $\times0.5$ to $\times2$\textsuperscript{a} & $0.00$\% & $0.00$\% \\
\midrule
Kept-edge ratio $r$, $\times0.5$   & $-0.77$\% [$-2.13$, $+0.58$] & $+0.16$\% [$-0.09$, $+0.41$] \\
Kept-edge ratio $r$, $\times0.75$  & $-0.66$\% [$-1.98$, $+0.65$] & $0.00$\% \\
Kept-edge ratio $r$, $\times1.5$   & $+0.23$\% [$-0.72$, $+1.17$] & $+0.08$\% [$-0.27$, $+0.43$] \\
Kept-edge ratio $r$, $\times2$     & $-0.74$\% [$-2.26$, $+0.78$] & $+1.08$\% [$+0.18$, $+1.98$] \\
\midrule
Community bound $\beta$, $\times0.5$  & $+1.77$\% [$-0.17$, $+3.72$] & $+0.63$\% [$-0.07$, $+1.32$] \\
Community bound $\beta$, $\times0.75$ & $+0.25$\% [$-1.19$, $+1.69$] & $+0.03$\% [$-0.46$, $+0.51$] \\
Community bound $\beta$, $\times1.5$  & $-0.25$\% [$-2.29$, $+1.79$] & $-0.45$\% [$-1.20$, $+0.30$] \\
Community bound $\beta$, $\times2$    & $-0.25$\% [$-1.60$, $+1.10$] & $+1.19$\% [$+0.25$, $+2.13$] \\
\midrule
All four, $\times0.5$ & $+1.15$\% [$-1.14$, $+3.43$] & $+0.59$\% [$-0.08$, $+1.27$] \\
All four, $\times2$   & $-1.10$\% [$-2.57$, $+0.38$] & $+1.78$\% [$+1.02$, $+2.54$] \\
\bottomrule
\end{tabularx}}
\tabnote{Entries are mean paired changes in station visits against the nominal setting, with 95\% confidence intervals; a positive value means more visits. Every setting returned a feasible layout on every instance. \textsuperscript{a}All four settings give bit-identical output, so the rows are collapsed into one.}
\label{tab:sensitivity}
\end{table}

\section{Paired statistical testing of the synthetic benchmark}\label{sup:tests}
Section~5 of the main article compares the methods on mean station visits, gaps, running time
and workload feasibility across the 29 synthetic instances. This section reports the paired
statistical testing behind those comparisons, which is kept out of the main text because its
resolution is bounded by the instance counts and adds little to the readings the main tables
already support.

Instances differ by two orders of magnitude in difficulty, so comparing two methods on their
averages would confound the between-method difference with the between-instance variance. We
therefore compare them instance by instance. For each pair of methods we take the per-instance
difference in station visits and test whether those differences are systematically signed. The
test is the two-sided exact Wilcoxon signed-rank test \citep{wilcoxon1945}, the standard choice
for comparing algorithms across a set of benchmark instances because it imposes no
distributional assumption on the differences \citep{demsar2006}, and ties are discarded. The
reported $p$ is the probability, under the null hypothesis that the paired differences are
symmetric about zero, of observing a signed-rank statistic at least as extreme as the one
obtained, in either direction. Table~\ref{tab:tests} reports the outcome against the
set-variable MILP, the per-instance reference of the main article.

\begin{table}[H]
\centering
\tbl{Paired Wilcoxon signed-rank tests on the 29 synthetic instances, against the set-variable MILP.}
{\small\begin{tabular*}{\textwidth}{@{\extracolsep{\fill}}llrrrr@{}}
\toprule
Comparison & Size (SKUs) & $n$ & Wins & $p$ & Min.\ $p$ \\
\midrule
Binary MILP vs.\ reference & 50      & 12 & 4  & 0.74   & 0.0005 \\
                           & 500     & 10 & 8  & 0.049  & 0.0020 \\
                           & 1{,}000 & 4  & 1  & 0.38   & 0.125  \\
                           & 2{,}000 & 3  & 0  & 0.25   & 0.25   \\
                           & pooled  & 29 & 13 & 0.91   & --     \\
\midrule
CG vs.\ reference          & 50      & 12 & 1  & 0.0010 & 0.0005 \\
                           & 500     & 10 & 2  & 0.0273 & 0.0020 \\
                           & 1{,}000 & 4  & 2  & 1.00   & 0.125  \\
                           & 2{,}000 & 3  & 3  & 0.25   & 0.25   \\
                           & pooled  & 29 & 8  & 0.18   & --     \\
\midrule
GA and SA-C vs.\ reference & 50      & 12 & 0  & 0.0005 & 0.0005 \\
                           & 500     & 10 & 0  & 0.0020 & 0.0020 \\
                           & 1{,}000 & 4  & 0  & 0.125  & 0.125  \\
                           & 2{,}000 & 3  & 0  & 0.25   & 0.25   \\
\midrule
Heuristic vs.\ reference   & 50      & 12 & 0  & 0.0005 & 0.0005 \\
                           & 500     & 10 & 0  & 0.0020 & 0.0020 \\
                           & 1{,}000 & 4  & 0  & 0.125  & 0.125  \\
                           & 2{,}000 & 3  & \textbf{3}  & 0.25   & 0.25   \\
\bottomrule
\end{tabular*}}
\tabnote{``Wins'' counts instances on which the first-named method returns fewer visits. ``Min.\ $p$'' is the smallest two-sided $p$ attainable at that sample size, $2^{1-n}$, and where $p$ equals it the test is saturated rather than failed, so the win count carries the evidence. The GA, SA-C and Heuristic rows coincide at 50, 500 and 1{,}000 SKUs because none of the three wins a single paired instance there. The Heuristic separates from them only at 2{,}000 SKUs, where it wins all three.}
\label{tab:tests}
\end{table}

The tests saturate at the two larger sizes. The resolution of the test is bounded by the number
of instances available, and on this benchmark that bound is binding: at 1{,}000 and 2{,}000
SKUs no result, however unanimous, can reach the conventional 0.05 threshold. At those two
sizes we therefore treat the win count as the primary evidence, which is the form in which the
main article states them.

The separations are reported without correction for multiplicity. Of the seven separations
below 0.05 that the table shows, the five at $p \le 0.0020$ survive a Bonferroni, Holm or
Benjamini--Hochberg correction over the eighteen tests reported, whereas the two at
$p = 0.0273$ and $p = 0.049$ do not. A small $p$ indicates that the methods differ, whereas a
large $p$ indicates only that these instances are insufficient to separate them, which is a
weaker statement than equivalence. Effect size and consistency are therefore reported
separately: the Gap column of the main article's benchmark table quantifies the magnitude of a
difference, and the test quantifies how reliably it recurs.

\section{Pseudocode of the greedy clustering heuristic}\label{sup:heuristic}
Algorithm~\ref{alg:heuristic} states in full the five stages that Section~4.1 of the main
article describes and its Figure~2 sketches. Here $|s|$ is the product count of station $s$,
$\mathrm{cnt}_{ij}$ the co-occurrence count of a product pair, $\mathrm{cnt}_i$ the order
frequency of product $i$, $\mathrm{load}(s)=\sum_{p\in s}L_p/V_s$ the workload of station $s$,
and $\zeta_s$, $T_s$, $L_p$ and $V_s$ are as defined in Section~3 of the main article. The four
thresholds are restated in the input line so that the procedure can be read on its own. The
ratio test in Step~1 normalises by the \emph{larger} of the two endpoint frequencies, so it is
symmetric in $i$ and $j$ and reads as the confidence of the weaker association direction.

\begin{algorithm}[H]
\caption{Greedy product clustering and station assignment}
\label{alg:heuristic}
\begin{algorithmic}[1]
\State \textbf{Input:} products $P$, orders $O$, stations $S$; filters from $N=|P|$: floor $\phi=\max(2,\lfloor N/100\rfloor)$, minimum co-occurrence $\gamma=\max(2,\lfloor N/50\rfloor)$, kept-edge ratio $r=0.1$; size bound from the slot capacity $\zeta=|P|/|S|$: $\beta=\max(5,\lfloor 0.4\,\zeta+\tfrac12\rfloor)$.
\Statex
\State \textbf{Step 1: preprocessing and pair generation}
\State $P^{\ast}\gets\{\,p\in P:\mathrm{freq}(p)\ge\phi\,\}$
\ForAll{multi-item orders $o\in O$}
  \State form every unordered pair $\{p_i,p_j\}\subseteq P_o\cap P^{\ast}$
\EndFor
\State for each pair, $\mathrm{cnt}_{ij}\gets\lvert\{\,o\in O:\{p_i,p_j\}\subseteq P_o\,\}\rvert$; keep pairs with $\mathrm{cnt}_{ij}\ge\gamma$
\State for each kept pair, $\eta_{ij}\gets\mathrm{cnt}_{ij}/\max(\mathrm{cnt}_i,\mathrm{cnt}_j)$; keep pairs with $\eta_{ij}\ge r$
\State sort the kept pairs by decreasing $\mathrm{cnt}_{ij}$
\Statex
\State \textbf{Step 2: community formation}
\State $\mathcal{C}\gets\varnothing$; mark every product unassigned
\While{candidate pairs remain}
  \State take the top pair $\{p_a,p_b\}$ and open a community $c\gets\{p_a,p_b\}$
  \While{$|c|<\beta$ and an unassigned product is adjacent to $c$}
    \State $p_x\gets$ unassigned product of highest aggregate co-occurrence with $c$
    \State $c\gets c\cup\{p_x\}$; delete every remaining pair that contains $p_x$
  \EndWhile
  \State $\mathcal{C}\gets\mathcal{C}\cup\{c\}$
\EndWhile
\Statex
\State \textbf{Step 3: post-processing}
\ForAll{products $p$ with $\mathrm{freq}(p)<\phi$}
  \If{some community $c$ has room ($|c|<\beta$)}
    \State add $p$ to the community of highest co-occurrence with $p$
  \EndIf
\EndFor
\State group any still-unassigned products into new communities of size at most $\beta$
\Statex
\State \textbf{Step 4: station assignment}
\State call $s$ \emph{admissible} for a block $c$ when $|s|+|c|\le\zeta_s$ and $\mathrm{load}(s)+\sum_{p\in c}L_p/V_s\le T_s$
\State order the communities by cumulative demand, most impactful first
\ForAll{communities $c$ in that order}
  \If{some station is admissible for $c$}
    \State assign $c$ to the admissible station of highest historical preference
    \State \Comment{ties broken by the lowest resulting load $\mathrm{load}(s)/T_s$}
  \Else
    \State split $c$, heaviest product first, placing each product by the same rule
  \EndIf
\EndFor
\Statex
\State \textbf{Step 5: workload repair} \Comment{$\mathrm{aff}(x,s)=\sum_{y\in s}\mathrm{cnt}_{xy}$}
\While{some station exceeds its workload cap}
  \State $s\gets$ the station of greatest relative overload $\mathrm{load}(s)/T_s$
  \State $E\gets$ exchanges of a heaviest product $p$ of $s$ against a lightest product $q$
  \State \hskip 1.2em of another station $t$ that lighten $s$ and keep $t$ inside its cap
  \State \textbf{if} $E=\varnothing$ \textbf{then break}
  \State apply the $(p,q)\in E$ of least broken weight $\mathrm{aff}(p,s)+\mathrm{aff}(q,t)-\mathrm{aff}(p,t)-\mathrm{aff}(q,s)$
\EndWhile
\State \textbf{Output:} a product-to-station assignment respecting $\zeta_s$ and $T_s$
\end{algorithmic}
\end{algorithm}

\bibliographystyle{apacite}
\bibliography{IJPR_CSLAP}

\end{document}
